% 埃米·诺特（综述）
% license CCBYSA3
% type Wiki

本文根据 CC-BY-SA 协议转载翻译自维基百科\href{https://en.wikipedia.org/wiki/Emmy_Noether}{相关文章}

\begin{figure}[ht]
\centering
\includegraphics[width=6cm]{./figures/c0d48f5d45fa6e41.png}
\caption{诺特，约摄于1900年至1910年间。} \label{fig_AMnt_1}
\end{figure}
阿玛莉·艾米·诺特（Amalie Emmy Noether，\(^\text{[a][b]}\)，1882年3月23日－1935年4月14日）是一位德国数学家，她在抽象代数领域作出了许多重要贡献。她还证明了诺特第一与第二定理，这两项成果在数学物理中具有基础性意义。\(^\text{[4]}\)帕维尔·亚历山德罗夫、阿尔伯特·爱因斯坦、让·迪厄多内、赫尔曼·外尔以及诺伯特·维纳都曾称她为“数学史上最重要的女性”。\(^\text{[5][6][7]}\)作为当时最杰出的数学家之一，诺特发展了环论、域论和代数理论。在物理学中，诺特定理揭示了对称性与守恒定律之间的深刻联系。\(^\text{[8]}\)

诺特出生于法兰克尼亚小镇埃尔朗根的一个犹太家庭，她的父亲是数学家马克斯·诺特。起初，她计划在通过所需考试后教授法语和英语，但后来转而在其父任教的埃尔朗根－纽伦堡大学学习数学。1907年，在保罗·戈尔丹的指导下，她完成了博士论文。此后，她在埃尔朗根数学研究所无薪工作了七年。\(^\text{[9]}\)当时，女性基本上被排除在学术职位之外。1915年，大卫·希尔伯特和费利克斯·克莱因邀请她加入哥廷根大学数学系——这是当时世界知名的数学研究中心。然而哲学学院对此提出反对，因此她在接下来的四年里只能以希尔伯特的名义讲课。1919年，她的任教资格最终获批，从而获得了“私人讲师”的职称。\(^\text{[9]}\)

诺特一直是哥廷根大学数学系的核心成员，直至1933年为止；她的学生们有时被称为“诺特男孩”。1924年，荷兰数学家B.L.范德瓦尔登加入她的学术圈，并很快成为她思想的主要阐述者；她的研究成果成为他1931年具有深远影响的教材《现代代数》第二卷的理论基础。到1932年在苏黎世举行的国际数学家大会上作大会报告时，诺特的代数才能已获得世界范围的认可。次年，德国纳粹政府将犹太人清除出大学职位，诺特被迫移居美国，在宾夕法尼亚州的布林茅尔学院任教。在那里，她教授研究生和博士后课程，学生包括玛丽·约翰娜·魏斯和奥尔加·陶斯基－托德。与此同时，她还在新泽西州普林斯顿的高等研究院讲学并从事研究工作。\(^\text{[9]}\)

诺特的数学研究通常被划分为三个“时期”。\(^\text{[10]}\)在第一个时期（1908–1919年），她对代数不变量理论和数域理论作出了重要贡献。她在变分法中的微分不变量研究——即著名的诺特定理——被称为“在指导现代物理学发展过程中最重要的数学定理之一”。\(^\text{[11]}\)在第二个时期（1920–1926年），她开始从事一项“改变了[抽象]代数面貌”的工作。\(^\text{[12]}\)在她1921年的经典论文《环域中的理想论》（Idealtheorie in Ringbereichen，意为“环域中理想的理论”）中，诺特将交换环中的理想理论发展成为一个具有广泛应用的数学工具。她巧妙地运用了升链条件，而满足该条件的数学对象也因纪念她而被称为“诺特环”。在第三个时期（1927–1935年），她发表了关于非交换代数与超复数的研究成果，并将群的表示论与模论和理想论统一起来。除了自己的论文之外，诺特还慷慨地分享思想，对多项由其他数学家发表的研究方向作出了重要贡献，甚至包括一些与她主要研究领域相距甚远的领域，如代数拓扑。
\subsection{传记}
\subsubsection{早年生活}
\begin{figure}[ht]
\centering
\includegraphics[width=6cm]{./figures/9bf717f586e25b88.png}
\caption{诺特在巴伐利亚的埃尔朗根市长大，图为1916年的一张该城市明信片。} \label{fig_AMnt_2}
\end{figure}
阿玛莉·艾米·诺特于1882年3月23日出生在巴伐利亚的埃尔朗根。\(^\text{[13]}\)她是数学家马克斯·诺特与伊达·阿玛莉亚·考夫曼的四个孩子中的长女，父母皆出身于富裕的犹太商人家庭。\(^\text{[14]}\)她的名字原本是“阿玛莉”，但从小便开始使用中间名“艾米”，并在成年后及其所有学术著作中始终沿用这一名字。\(^\text{[b]}\)

年轻时的诺特在学业上并不特别突出，但以聪慧与友善著称。她有些近视，童年时期说话带轻微的口齿不清。多年后，一位家庭朋友回忆说，在一次儿童聚会上，小诺特曾迅速解开一道益智谜题，显示出她早期出众的逻辑才能。\(^\text{[15]}\)她像当时大多数女孩一样学习烹饪与家务，也上过钢琴课。虽然她对这些活动都缺乏热情，但却非常热爱舞蹈。\(^\text{[16]}\)
\begin{figure}[ht]
\centering
\includegraphics[width=6cm]{./figures/40e1cdf2d55a10d2.png}
\caption{艾米·诺特与她的三位弟弟——阿尔弗雷德、弗里茨和罗伯特，摄于1918年之前。} \label{fig_AMnt_3}
\end{figure}
诺特有三个弟弟。长弟阿尔弗雷德·诺特生于1883年，1909年在埃尔朗根获得化学博士学位，但九年后便去世。\(^\text{[17]}\)次弟弗里茨·诺特生于1884年，在慕尼黑求学，并在应用数学领域作出过贡献。\(^\text{[18]}\)据推测，他于第二次世界大战期间的1941年在苏联遭到处决。\(^\text{[19]}\)最小的弟弟古斯塔夫·罗伯特·诺特生于1889年，关于他的生平所知甚少；他长期患有慢性疾病，并于1928年去世。\(^\text{[20][21]}\)
\subsubsection{教育经历}
诺特自幼在法语与英语方面展现出优秀的天赋。1900年初，她参加了外语教师资格考试，并获得了“Sehr gut”（非常好）的总评成绩。这一成绩使她有资格在女子学校教授外语，但她选择继续在父亲任教的埃尔朗根－纽伦堡大学深造数学。[22][23]

这一决定在当时极不寻常。就在两年前，该校学术参议院曾宣称，允许男女同校将会“颠覆整个学术秩序”。\(^\text{[24]}\)在这所拥有986名学生的大学里，诺特是仅有的两位女性之一。她只能旁听课程，不能正式注册听课，并且每听一门课都必须得到任课教授的特别许可。尽管困难重重，她仍于1903年7月14日通过了纽伦堡实科中学的毕业考试。\(^\text{[22][25][26]}\)

在1903至1904年的冬季学期中，她前往哥廷根大学学习，听取了天文学家卡尔·施瓦西以及数学家赫尔曼·闵可夫斯基、奥托·布卢门塔尔、费利克斯·克莱因和大卫·希尔伯特的讲座。\(^\text{[27]}\)
\begin{figure}[ht]
\centering
\includegraphics[width=6cm]{./figures/8f7a5161e1329585.png}
\caption{保罗·戈尔丹指导了诺特的博士论文，研究主题为四次型不变量。} \label{fig_AMnt_4}
\end{figure}
1903年，巴伐利亚大学对女性全面入学的限制被取消。\(^\text{[28]}\)诺特返回埃尔朗根，并于1904年10月正式重新入学，声明她将专注于数学研究。当年她所在年级共有六名女性（其中两名为旁听生），而她是唯一一位选择数学专业的女性。\(^\text{[29]}\)在保罗·戈尔丹的指导下，她于1907年完成了博士论文《三元四次型不变量系的构成》，并于同年以“最优等成绩”毕业。\(^\text{[30][31]}\)戈尔丹属于“不变量计算学派”的代表人物，而诺特的论文以列出三百多个明确计算出的不变量作为结尾。这种具体计算的方法后来被希尔伯特所开创的更抽象、更一般化的不变量理论方法所取代。\(^\text{[32][33]}\)尽管她的论文在当时受到好评，但诺特后来却将这篇论文以及她之后撰写的几篇类似论文称为“垃圾”。\(^\text{[33]}\)她此后的全部研究工作都进入了一个完全不同的领域。\(^\text{[34]}\)
